%%======================================================================
\section{What Healthcare Needs from Humanoid Robots}
\label{sec:applications}
%%======================================================================

Each clinical role imposes distinct requirements on a humanoid platform. The
clinical studies identified in our search are evidence of how far current systems
come toward meeting those requirements, not the subject of this section in
themselves. We therefore ask not ``what has been done?'' but ``how close is the
field to what clinical deployment actually demands?'' The gap between those two
positions is a central concern of this review.

\subsection{Clinical Procedures and Surgical Assistance}
\label{sec:procedures}

Clinical procedure assistance demands sub-millimeter manipulation precision,
controlled interaction forces on the order of 1--2~N for delicate soft-tissue
manipulation~\cite{taylor2003medical,okamura2009haptic}, sterility compliance
(no bare actuator surface in the sterile field)~\cite{taylor2003medical}, and
consistent task execution over procedure durations of 20--120~minutes, 
consistent with the range of procedures evaluated in this  
review~\cite{atar2025humanoids,liang2025lapsurgie,cho2026humanoid}. For
teleoperated roles, latency below 100~ms round-trip is required to avoid  
operator disorientation during fine tool 
manipulation~\cite{lum2009latency,avgousti2016medical,motiwala2025telesurgery}. 
For autonomous or semi-autonomous roles, real-time anatomical landmark 
detection and adaptive trajectory generation under tissue deformation are 
prerequisites~\cite{yang2017medical,saeidi2022autonomous}.

Current evidence documents how far the field has advanced toward these
requirements, and where the gaps remain. Three studies, all using the Unitree
G1 under teleoperation, demonstrate feasibility across a range of procedures:
Atar et al.~\cite{atar2025humanoids} achieved approximately 70\% task success
across seven clinical procedures (auscultation, ultrasound, intravenous
catheter insertion, and four others) with non-clinician teleoperators; Liang et
al.~\cite{liang2025lapsurgie} demonstrated laparoscopic instrument manipulation
with inverse-kinematics remapping and enforced remote-center-of-motion
constraint; and Cho et al.~\cite{cho2026humanoid} performed cadaveric
endoscopic surgery with a humanoid as first assistant. The robot's contribution
in each case is physical presence and dexterity; clinical judgment and
real-time adaptation are supplied by the remote operator. Tasks requiring high
force precision (catheter insertion, finger-stick testing) showed the highest
failure rates, attributable to actuator force resolution limits and insufficient
tactile feedback.

The requirements gap is precise. Teleoperation architecture meets the
near-term procedure assistance requirement, but the 70\% task success rate is
insufficient for unassisted clinical use. No study has measured forces in
tissue contact against published safety thresholds. All evidence is pre-clinical
(mannequins, cadavers, test fixtures); no patient data exist. Independent
replication on any platform other than the Unitree G1 has not occurred.
\textbf{TRL: 3--4.}

\subsection{Elderly and Nursing Care}
\label{sec:eldercare}

Nursing care imposes a demanding combination of physical and social requirements.
On the physical side: safe patient handling (transfer, repositioning, ambulation
assist) requires payload capacity of 30--80~kg with compliant force
control~\cite{mukai2010riba}; gait stability on wet or uneven ward surfaces is
required~\cite{niosh2011slips,joeoh2019terrain}; and proximity to
patients with fragile physiology demands contact-force limits more restrictive
than industrial safety standards~\cite{iso13482,iso10218_1}. On the social side: sustained interaction
quality (naturalness, responsiveness, memory across sessions) matters for
patient trust and care acceptance~\cite{zafrani2022betweenfear}.

Against these requirements, the evidence reveals a clear bottleneck. The only
bipedal humanoid designed explicitly for nursing is RHP
Friends~\cite{benallegue2025rhp}, which demonstrates autonomous-to-teleoperated
handoff and nursing-specific motion primitives, but in controlled exhibition
conditions, not patient-facing deployment. The SPRING project provides the only
large-scale eldercare deployment data (60+ participants over 4 years), but on a
wheeled ARI platform~\cite{alamedapineda2024spring}; longer-term telepresence
robot deployments in eldercare settings have similarly relied on wheeled
systems~\cite{cesta2016amicare}. Whether an upright humanoid form factor confers
a user acceptance advantage in this setting remains a hypothesis that bipedal
systems have not yet verified with their own evidence. Imtiaz and
Khan~\cite{imtiaz2024perceptions} surveyed 269 US nursing
home administrators and found the primary deployment barriers are economic
(cost, lack of effectiveness data), not technical; Mitzner et
al.~\cite{mitzner2018technology} similarly found that older adults' concerns
center on capability and integration rather than technology rejection, a finding
that reframes the investment priority: the field must produce health-economic
evidence, not only engineering demonstrations, yet rigorous economic evaluations
even for existing aged-care technologies remain
scarce~\cite{bulamu2018economic}. \textbf{TRL: 4 for demonstrated hardware;
TRL~4--5 for wheeled platform deployment evidence (scope-constrained).}

\subsection{Rehabilitation}
\label{sec:rehabilitation}

Rehabilitation demands contact safety as a non-negotiable prerequisite: no
patient-contact study on a full-size bipedal humanoid is ethically or
regulatorily feasible without a validated compliance framework that bounds peak
interaction forces, the requirement codified for rehabilitation robots in
IEC~80601-2-78 and, for collaborative force limits, in
ISO/TS~15066~\cite{iec80601_2_78,isots15066,yang2017medical}. Beyond safety, repetition consistency across hundreds of
movement cycles (neuromotor rehabilitation protocols target repetition doses on
the order of 300+ per session, far above the $\sim$30 typically delivered in
conventional therapy)~\cite{lang2009observation,birkenmeier2010translating},
patient-adaptive trajectory modification in response to
effort and fatigue, and integration with standardized outcome measures
(Fugl-Meyer Assessment, 10-meter walk
test)~\cite{fuglmeyer1975poststroke,winstein2016guidelines} define the clinical
requirements for this role. The evidence bar for robotic rehabilitation is well established:
Lo et al.~\cite{lo2010robot} demonstrated in a multi-site randomized trial that
robot-assisted therapy for post-stroke upper-limb impairment can achieve
outcomes comparable to standard care; the MIT-Manus neurorehabilitation
robot~\cite{krebs1998robot} and ARMin exoskeleton RCT of
Klamroth-Marganska et al.~\cite{klamroth2014armin} similarly demonstrated
TRL~7--8 performance for upper-limb rehabilitation. Robotic rehabilitation
technology reviews confirm this trajectory~\cite{poli2013robotic}. These serve
as the reference standard that bipedal humanoid rehabilitation systems must
eventually match.

The requirements gap here is in two layers. The first is the safety precondition.
Lu et al.~\cite{lu2025gentlehumanoid} address it directly, demonstrating
measurable peak force reduction during quasi-static upper-body contact under the
GentleHumanoid spring compliance framework on the Unitree G1. This is the
necessary engineering precondition for all patient-contact rehabilitation work
that follows, but it has not been validated in dynamic motion or replicated on
other platforms. The second layer is evidence of therapeutic efficacy. Nguyen et
al.~\cite{nguyen2024eeg} establish a neuroscientific rationale: humanoid
observation elicits sensorimotor EEG responses in healthy subjects consistent with
action observation therapy (AOT) for stroke rehabilitation. Sobrepera et
al.~\cite{sobrepera2022feasibility} ($n\,=\,6$) show that patients preferred
humanoid telepresence over video for neuromotor rehabilitation, attributing this
to physical embodiment. The evidence is pre-clinical throughout: the EEG study
enrolled healthy subjects only; no study applies a validated clinical outcome
measure; no study has involved patient contact with a current-generation full-size
bipedal platform. \textbf{TRL: 3--4.}

\subsection{Mental Health and Social Interaction}
\label{sec:mentalhealth}

Mental health support demands naturalness in unscripted conversation, consistent
interaction quality across sessions, and behavioral adaptation to individual
patient affect and disclosure patterns~\cite{fitzpatrick2017woebot}. Trust is a functional requirement, not a
preference: a patient who distrusts the system will not disclose, making the
interaction therapeutically inert~\cite{lucas2014computer}. For clinical deployment, outcomes must be
measurable with validated instruments
(PHQ-9~\cite{kroenke2001phq9}, GAD-7~\cite{spitzer2006gad7},
PANSS~\cite{kay1987panss}) in clinically diagnosed populations.

The evidence presents an inverse quality-specificity tradeoff. The
best-designed studies are randomized: Robinson et al.~\cite{robinson2023brief}
report a pilot RCT ($n\,=\,230$) in which an autonomous humanoid delivered a
mindful breathing intervention with positive wellbeing outcomes; Sayis and
Gunes~\cite{sayis2024technology} find that only the humanoid robot condition
produced statistically significant mood improvements and increased self-disclosure
($n\,=\,42$). Neither study specifies the humanoid platform, which prevents
extrapolation to full-size bipedal systems. The only full-size platform study
(Ghosh et al.~\cite{ghosh2024usercentered}, $n\,=\,14$) uses NAO (58~cm, out of
scope for physical care roles). All studies use subclinical or university wellness
populations; none involves a clinically diagnosed population under treatment. The
embodiment advantage is consistent across designs, but it has been demonstrated
on unspecified or out-of-scope platforms, in non-clinical settings, with
non-patients. Yuan et al.~\cite{yuan2025integrating} provide a transferable LLM
personalization framework for dementia care on Pepper, and Lindsay et
al.~\cite{lindsay2022socially} demonstrate automated clinical planning in
pediatric settings; both represent relevant methodological precedent for bipedal
deployment. \textbf{TRL: 4--5 for social wellbeing applications in non-clinical
settings; TRL~2--3 for bipedal humanoid deployment in clinical mental health
contexts.}

\subsection{Hospital Logistics}
\label{sec:logistics}

Hospital logistics requires reliable wayfinding in dynamic ward environments,
payload capacity for transport tasks (medication carts, linen, supply bins up to
20--30~kg), and compatibility with existing hospital information systems for task
scheduling~\cite{ozkil2009service,bloss2011mobile}. Crucially, it does not require bimanual dexterity, standing height,
or natural language sophistication, which raises the threshold question: is the
humanoid form factor the right choice for this role at all?

Our systematic search (Section~\ref{sec:intro}) identified no qualifying papers
for bipedal humanoid hospital logistics in the 2022--2026 period
(Figure~\ref{fig:evidence_dist}). Wheeled autonomous mobile robots already
address this role with operational reliability proven over three decades of
hospital deployment, from the HelpMate courier of the 1990s to current
delivery fleets and sterile instrument logistics
systems~\cite{evans1994helpmate,bloss2011mobile,fragapane2021autonomous}; they
provide no demonstrated functional disadvantage for logistics tasks, and their
simpler regulatory profile (autonomous ground vehicle versus medical device)
means faster deployment.
The humanoid form factor offers no structural advantage for corridor navigation,
supply transport, or medication delivery over wheeled alternatives. A study of cobot logistics in a 900-bed hospital system found that nurses
report primarily workflow integration concerns rather than form-factor
preferences~\cite{birkhoff2024integrating}. Academic research attention in
2022--2026 was directed toward clinical procedure and HRI domains, reflecting this assessment implicitly. \textbf{TRL: 2 for bipedal
humanoid hospital logistics; the role is better addressed by non-humanoid
platforms at current maturity.}

\subsection{Requirements Gap Analysis}
\label{sec:crossdomain}

Table~\ref{tab:evidence} summarizes the requirements-versus-evidence position
across all five domains.

\begin{table*}[t]
  \captionsetup{justification=centering,skip=2pt}
  \caption{Requirements Gap by Clinical Domain (March 2026).}
  \label{tab:evidence}
  \setlength{\tabcolsep}{3pt}
  \resizebox{\linewidth}{!}{%
  \begin{tabular}{@{}>{\raggedright\arraybackslash}p{2.95cm}
                  >{\raggedright\arraybackslash}p{7.70cm}
                  >{\raggedright\arraybackslash}p{4.25cm}
                  >{\raggedright\arraybackslash}p{4.85cm}
                  >{\raggedleft\arraybackslash}p{2.20cm}
                  >{\raggedleft\arraybackslash}p{1.20cm}@{}}
    \toprule
    \textbf{Domain} & \textbf{Key Papers} & \textbf{Platform} & \textbf{Study Type} & \textbf{N} & \textbf{TRL} \\
    \midrule
    Clinical procedures &
      Atar et al.~\cite{atar2025humanoids}, Liang et al.~\cite{liang2025lapsurgie}, Cho et al.~\cite{cho2026humanoid} &
      Unitree G1 & Lab; user study; cadaver & 0 patients & 3--4 \\[3pt]
    Elderly / nursing &
      Benallegue et al.~\cite{benallegue2025rhp}, Alameda-Pineda et al.\textsuperscript{*}~\cite{alamedapineda2024spring} &
      RHP Friends; ARI\textsuperscript{*} & Exhibition demo; field deploy &
      0 (demo); 60+\textsuperscript{*} & 4; 4--5\textsuperscript{*} \\[3pt]
    Rehabilitation &
      Nguyen et al.~\cite{nguyen2024eeg}, Sobrepera et al.~\cite{sobrepera2022feasibility}, Lu et al.~\cite{lu2025gentlehumanoid} &
      Humanoid; SRAT; Unitree G1 & EEG pilot; case series; engineering &
      0 / 6 / 0 & 3--4 \\[3pt]
    Mental health / HRI &
      Robinson et al.~\cite{robinson2023brief}, Sayis and Gunes~\cite{sayis2024technology} &
      Unspecified humanoid & RCT; comparative & 230; 42 & 4--5 \\[3pt]
    Hospital logistics & \multicolumn{4}{l}{\emph{No qualifying papers; wheeled systems dominate (TRL~2)}} & 2 \\
    \bottomrule
    \multicolumn{6}{@{}l@{}}{\textbf{TRL}: ISO 16290:2013 scale.} \\
    \multicolumn{6}{@{}l@{}}{1 = basic principles; 2 = concept formulated; 3 = experimental PoC; 4 = lab validated; 5 = relevant-environment validated; 7 = operational prototype; 9 = proven.} \\
    \multicolumn{6}{@{}l@{}}{\textsuperscript{*}ARI/SPRING is a wheeled platform, not bipedal; scope caveat applies.} \\
  \end{tabular}%
  }
\end{table*}

Three structural patterns constrain the generalizability of the entire evidence
base. \textbf{Platform monoculture:} the Unitree G1 is the experimental platform
in four of the clinical studies (all three procedure studies and the compliance
rehabilitation study). No other commercial platform appears in more than one qualifying study,
because the Unitree G1's open SDK is currently the only commercial research
access point; until equivalent access exists for higher-dexterity platforms
(Sanctuary Phoenix Gen~8, 1X NEO Beta), cross-platform replication is
structurally blocked. \textbf{Geographic concentration:} studies originate
predominantly from the USA (8 papers), Japan (2), and the UK (2); no study from
continental Europe (excluding the SPRING program), the Global South, or Southeast
Asia appears in the qualifying set. Healthcare delivery models, regulatory
frameworks, and workforce demographics differ substantially across these absent
geographies. \textbf{Methodological ceiling:} no qualifying study constitutes a
randomized controlled trial with a bipedal humanoid in a clinical population; no
study applies a validated clinical outcome measure as its primary endpoint; no
study exceeds a single session; no study includes health economic data. These
gaps collectively define the minimum evidence package for any future regulatory
submission: a prospective, multi-session, clinically powered study in a
patient-contact context, with pre-registered primary endpoints and a costing
analysis.

The domain-level TRL picture confirms that clinical procedures and rehabilitation
share the same bottleneck (contact safety and force control), while elderly
care and mental health share a different bottleneck: the bottleneck is not
hardware but evidence. Both domains have transferred results from non-bipedal
platforms (ARI/SPRING, unspecified humanoids) and non-clinical populations that
must be replicated on full-size bipedal systems in patient populations before
regulatory evidence packages can be assembled.

Figure~\ref{fig:evidence_dist} visualizes the distribution of the qualifying
clinical studies by domain and year, with marker shape indicating study design type.

\begin{figure*}[t]
  \centering
  \Description{Dot matrix showing 14 qualifying papers organized by clinical
  domain (y-axis: Clinical Procedures, Elderly/Nursing Care, Rehabilitation,
  Mental Health/HRI) and publication year (x-axis: 2022-2026). Marker shape
  encodes study design: circle for engineering/lab, diamond for case series or
  observational, star for RCT/comparative. Clinical procedure papers are
  concentrated in 2025-2026; mental health studies span 2023-2025;
  rehabilitation studies span 2022-2025; elderly/nursing care studies
  span 2024-2025. Author surnames label each marker.}
  \includegraphics[width=\textwidth]{figures/evidence_landscape.pdf}
  \caption{Evidence distribution: qualifying clinical studies by clinical domain and
  publication year. Marker shape encodes study design (circle: engineering/lab;
  diamond: case series/observational; star: RCT/comparative); color encodes
  domain. Clinical procedure evidence is concentrated in 2025--2026.
  Rehabilitation and mental health evidence is more temporally distributed
  but limited by platform specificity. No papers qualify for hospital
  logistics.}
  \label{fig:evidence_dist}
\end{figure*}
